\newif\ifglobal

%%%%%%%%%%%%%%%%%%%%%%%
%%% Compile to view %%%
%%%%%%%%%%%%%%%%%%%%%%%

%%%%%%%%%%%%%%%%%%%%%%%%%%%%%%%%%%%%%%%%%%%%%%%%%%%
%%% Readme:                                     %%%
%%% https://www.overleaf.com/read/bwdjvfjksvjy/ %%%
%%%%%%%%%%%%%%%%%%%%%%%%%%%%%%%%%%%%%%%%%%%%%%%%%%%

% >>> M?: marks a module setting number or important notice

% >>> M4: Set {./relative-path} to external files for this module <<<
\def \modulefiles {./28-HMAC}
% To add an external file, use:
% (Replace '---' with actual filename)
% '\input{\modulefiles/tables/---.tex}' for tables
% '\includegraphics{\modulefiles/figures/---}' for figures
% '\subfile{\modulefiles/---}' for register files

% >>> M1: This module is in global project? <<<
% 'YES' - uncomment; 'NO' - comment out
\globaltrue

\ifglobal
    \documentclass[main\_\_CN.tex]{subfiles}

\else
    % Font "FandolSong-Regular" used by ctex does not have GSUB table.
    % It causes warnings that do not affect compilation and can be ignored.
    \PassOptionsToPackage{quiet}{fontspec}

    \documentclass[openany, 10pt]{book}

    % >>> M2: In ./chip-spec-settings/preamble-trm-module, also find
    %         settings for visibility of todo notes, watermark <<<
    \usepackage{preamble-trm-module}


    % Enable Chinese typesetting
    \usepackage{ctex}

    % >>> M3: Temporary labels in Chinese
    %         you can add your own labels there <<<
    \myexternaldocument{temp-labels-cn}

\fi %\ifglobal


\setcounter{secnumdepth}{4}

% >>> M5: If this module needs any updates in
% './chip-spec-settings/preamble-trm-module.sty'
% (include additional package, etc.), contact your module PM
% or create the file './chip-spec-settings/changelog.tex' make the
% necessary updates yourself and document them in this file <<<

\raggedright


\begin{document}


% Sets language into which pre-defined bits of text are translated
\selectlanguage{chinese}

% Do not delete this block
\ifglobal
% Add the following front matter if
% this module is outside of global project
\else
    \InputIfFileExists{readme}{}{}
    {\let\clearpage\relax \listoftodos}
    {\let\clearpage\relax \listofdonetodos}
    \tableofcontents
    %\listoftables
    %\listoffigures
\fi


%%%%%%%%%%%%%%%%%%%%%%%%%
%%% TRM Module Begins %%%
%%%%%%%%%%%%%%%%%%%%%%%%%

\chapter{HMAC 加速器 (HMAC)}
\label{mod:hmac}
\hypertarget{hmac}{}

HMAC 模块通过 hash 算法 SHA-256 和密钥计算得到数据信息的信息认证码 (MAC)，采用 RFC 2104 标准。长度为 256 位的 HMAC 密钥存储在 eFuse 的密钥块中，可配置为读保护，即用户无法直接从 HMAC 加速器外部访问密钥。

\section{主要特性}
\begin{itemize}
    \item 使用标准 HMAC-SHA-256 算法
    \item 基于 eFuse 密钥的 HMAC-SHA-256 算法计算，
    \begin{itemize}
    \item 为提高安全性，软件在下行模式下无法访问其计算结果
    \item 软件在上行模式下可访问其计算结果
    \end{itemize}
    \item 下行模式下，支持生成 RSA 数字签名外设 (RSA\_DS) 所需的密钥
    \item 下行模式下，支持重启软禁用的 JTAG
\end{itemize}


\section{功能描述}
HMAC 模块有两种工作模式，即上行模式和下行模式。上行模式中，由用户提供 HMAC 信息并回读计算结果；下行模式中，HMAC 模块为其他两个可能的内部硬件模块提供输入：一方面，如 JTAG 被暂时禁用，HMAC 可使能 JTAG；另一方面，HMAC 可用作存储在 RSA\_DS 外设 flash 中数字签名参数的解密密钥。此外，这些计算在下行模式下自动进行，只要配置正确，就可以确保密钥及派生密钥材料的保密性。
\begin{tiplisting}
芯片释放复位信号后，HMAC 模块将检查 eFuse 中是否烧写有 RSA\_DS 外设所需要的功能密钥。如果存在该密钥，HMAC 模块将自动进入下行 RSA\_DS 模式，完成相应密钥计算。
\end{tiplisting}

\subsection{上行模式}

在上行模式下计算 HMAC 值，用户应执行以下步骤：

\begin{enumerate}
    \item 初始化 HMAC 模块，进入上行模式。
    \item 将正确填充的信息写入 HMAC 中，一次写入一个块。
    \item 从 HMAC 回读计算结果。
\end{enumerate}

有关此过程的详细步骤，可参见章节 \ref{sec:hmac-process-detailed}。

\begin{tiplisting}
使用上行模式的较多为支持 HMAC-SHA-256 算法的挑战-应答协议。假设挑战-应答协议中的通讯双方分别称为 A 和 B，想要交换的数据信息为 M，则协议的一般验证流程为：

\begin{itemize}
    \item A 计算出一个特殊的随机数 M。
    \item A 将 M 发送给 B。
    \item B（通过 M 和密钥）计算 HMAC 结果，并将结果发送给 A。
    \item A（通过 M 和密钥）内部计算 HMAC 结果。
    \item A 比较两次计算结果。如结果相同，则验证通过 B 的身份。
\end{itemize}
\end{tiplisting}

\subsection{下行模式－JTAG 启动功能}

JTAG 可被 eFuse 暂时关闭，关闭后，用户可使用 HMAC 模块重启 JTAG 功能。（详细步骤可参见章节 \ref{mod:efuse} \textit{\nameref{mod:efuse}}。）该模式下，用户应提供某 eFuse 密钥的 HMAC 计算结果，HMAC 模块将检查用户提供的数值与计算的数值是否匹配。如果二者相同，JTAG 将被启动，直到用户再次调用 HMAC 模块清除计算结果并关闭 JTAG。

用户可执行以下步骤重启 JTAG：

\begin{enumerate}
    \item 通过初始化时钟和 HMAC 复位信号启动 HMAC 模块，并通过配置 \hyperref[regdesc:HMACSETPARAPURPOSEREG]{HMAC\_SET\_PARA\_PURPOSE\_REG} 位进入下行 JTAG 启动模式，等待计算完成。详细步骤可参见章节 \ref{sec:hmac-process-detailed}。
    \item 将 1 写入 \hyperref[regdesc:HMACSOFTJTAGCTRLREG]{HMAC\_SOFT\_JTAG\_CTRL\_REG} 寄存器，进入 JTAG 重启模式。
    \item 将 256 位的 HMAC 值写入寄存器 \hyperref[regdesc:HMACWRJTAGREG]{HMAC\_WR\_JTAG\_REG}。该值是预先在本地使用 SHA-256 和已被写入 eFuse 的密钥对 32 字节的 0x00 进行 HMAC 计算得到的。需以大端字序分 8 次写入，每次 32 位。
    \item 如果由 eFuse 中密钥计算的 HMAC 结果与用户在步骤 3 中写入的数值一致，则 JTAG 重启。否则，JTAG 保持关闭状态。
    \item 在用户将 1 写入寄存器 \hyperref[regdesc:HMACSETINVALIDATEJTAGREG]{HMAC\_SET\_INVALIDATE\_JTAG\_REG} 或复位重启之前，JTAG 将保持关闭状态。如需再次重启 JTAG，请重复上述步骤。
    \end{enumerate}


\subsection{下行模式－RSA\_DS 密钥派生功能}

RSA 数字签名外设 (RSA\_DS) 使用 AES-CBC 算法加密其参数。HMAC 模块作为密钥导出函数 (KDF) 导出解密上述参数的 AES 密钥。

在启用 RSA\_DS 外设之前，用户必须先通过 HMAC 模块计算得到 RSA\_DS 外设工作时所需的密钥。详细信息可参见章节 \ref{mod:digsig} \textit{\nameref{mod:digsig}}。芯片上电后，HMAC 模块将检查 eFuse 中是否烧写有 RSA\_DS 外设所需要的功能密钥。如果存在该密钥，HMAC 模块将自动进入下行 RSA\_DS 模式，完成相应密钥计算。

\subsection{HMAC eFuse 配置}
\label{sec:hmac-config-function}

每个烧写入 eFuse 密钥块的 HMAC 密钥都有其功能，指定该密钥用于实现哪一功能。HMAC 中，未配置匹配的功能数值的密钥无法执行计算。当前 HMAC 模块共支持三种功能：下行模式下的 JTAG 重启功能和 RSA\_DS 密钥导出功能，以及上行模式下的 HMAC 计算功能。表 \ref{tab:hmac-key-purpose} 列出了上述各项功能，以及一种可指定密钥同时用于 JTAG 重启和 RSA\_DS 密钥导出的功能。

启动 HMAC 模块进行计算之前，用户应先读取章节 \ref{mod:efuse} \textit{\nameref{mod:efuse}} 中 EFUSE\_KEY\_PURPOSE\_x 的值（基于 eFuse 中共存在 6 个密钥，故 x 的值为 0 \verb+~+ 5，其中 \hyperref[fielddesc:EFUSEKEYPURPOSE0]{EFUSE\_KEY\_PURPOSE\_0} \verb+~+ \hyperref[fielddesc:EFUSEKEYPURPOSE1]{EFUSE\_KEY\_PURPOSE\_1} 属于寄存器 \hyperref[regdesc:EFUSERDREPEATDATA1REG]{EFUSE\_RD\_REPEAT\_DATA1\_REG}，\hyperref[fielddesc:EFUSEKEYPURPOSE2]{EFUSE\_KEY\_PURPOSE\_2} \verb+~+ \hyperref[fielddesc:EFUSEKEYPURPOSE5]{EFUSE\_KEY\_PURPOSE\_5} 属于寄存器 \hyperref[regdesc:EFUSERDREPEATDATA2REG]{EFUSE\_RD\_REPEAT\_DATA2\_REG}），检查 eFuse 中是否已烧写密钥。例如，在上行模式下，如果 EFUSE\_KEY\_PURPOSE\_0\verb+~+5 中没有 EFUSE\_KEY\_PURPOSE\_HMAC\_UP，则说明 eFuse 中没有上行模式可用的密钥。此时，可按照下述步骤将密钥烧写入 eFuse 中：

\begin{enumerate}
    \item 准备一个 256 位 HMAC 密钥，将其烧写到空的 eFuse 密钥块 $y$ 中。eFuse 中共有 6 个密钥块，分别为 4 \verb+~+ 9，因此 $y$ 的值为 4 \verb+~+ 9。也就是说，若密钥为 key0，则对应的密钥块为 block4。接着，将功能对应数值烧写至 EFUSE\_KEY\_PURPOSE\_$(y-4)$。例如在上行模式下，烧写密钥后，应将 EFUSE\_KEY\_PURPOSE\_HMAC\_UP（对应值为 8）写入 EFUSE\_KEY\_PURPOSE\_$(y-4)$。更多烧写 eFuse 密钥的详细信息，可参见章节 \ref{mod:efuse} \textit{\nameref{mod:efuse}}。
    \item 如有需要，配置 eFuse 密钥块读保护功能，使用户无法读取密钥值。用户应保留步骤中生成的随机密钥副本以便验证。
\end{enumerate}

请注意，功能为 EFUSE\_KEY\_PURPOSE\_HMAC\_DOWN\_ALL 的密钥既可用于 JTAG 重启，也可用于 RSA\_DS 密钥导出。

\input{\modulefiles/tables/HMAC-key-purpose__CN}

\textbf{选取 eFuse 的密钥块与 HMAC 功能}
\label{sec:hmac-config-purposes}

eFuse 控制器共有 6 个密钥块，即 KEY0 \verb+~+ 5 。用户将编号 \textcolor{red}{n} 写入寄存器 \hyperref[regdesc:HMACSETPARAKEYREG]{HMAC\_SET\_PARA\_KEY\_REG}，表示选择 KEY\textcolor{red}{n} 作为本次 HMAC 模块运行时使用的密钥。

将所使用功能对应的数值写入寄存器 \hyperref[regdesc:HMACSETPARAPURPOSEREG]{HMAC\_SET\_PARA\_PURPOSE\_REG} 完成配置（可参见章节 \ref{sec:hmac-process-detailed}）。需要注意的是，eFuse 存储器中的密钥在烧写时都定义了功能用途，只有当 HMAC 的配置功能与 KEY\textcolor{red}{n} 定义的功能用途相匹配时，HMAC 模块才会执行配置好的计算任务。否则，HMAC 模块将返回匹配错误结果并结束当前计算任务。

比如，如果用户选择了 KEY3 作为本次计算的密钥，且烧写入 \hyperref[fielddesc:EFUSEKEYPURPOSE3]{EFUSE\_KEY\_PURPOSE\_3} 中的数值为 6，即功能 EFUSE\_KEY\_PURPOSE\_HMAC\_DOWN\_JTAG。参照表 \ref{tab:hmac-key-purpose} 可知，KEY3 是用于 JTAG 重启的密钥。如果此时寄存器 \hyperref[regdesc:HMACSETPARAPURPOSEREG]{HMAC\_SET\_PARA\_PURPOSE\_REG} 中配置的数值也是 6，HMAC 模块才会启动 JTAG 重启功能的计算。


\tagged{ESP32-P4-latest}{\textbf{从密钥管理器中选取密钥}
\label{sec:hmac-config-km-key}

用户可以在密钥管理器中部署 HMAC 的密钥，部署方式请参考章节 \ref{mod:keymng} \textit{\nameref{mod:keymng}}。要选择密钥管理器中的该密钥，请将密钥编号 \textcolor{red}{7} 写入注册表 \hyperref[regdesc:HMACSETPARAKEYREG]{HMAC\_SET\_PARA\_KEY\_REG}。用户从密钥管理器中选取密钥时，注册表中 HMAC 操作的预定目标会自动设置为 ``HMAC Calculation"，其他设置都不会生效。
}

\subsection{调用 HMAC 流程（详细说明）}
\label{sec:hmac-process-detailed}
用户调用 \chipname{} 中 HMAC 流程如下：

\subsubsection{启动 HMAC 模块}
\label{sec:hmac-enable-module}

    \begin{enumerate}
        \item 置位寄存器 \hyperref[fielddesc:PCRHMACCLKEN]{PCR\_HMAC\_CLK\_EN} 中 HMAC 和 SHA 的外设时钟位，清除寄存器 \hyperref[fielddesc:PCRHMACRSTEN]{PCR\_HMAC\_RST\_EN} 中相应的外设复位位。关于如何配置时钟和重启信号，请参见章节 \ref{mod:resclk} \textit{\nameref{mod:resclk}}。
        \item 将 1 写入寄存器 \hyperref[regdesc:HMACSETSTARTREG]{HMAC\_SET\_START\_REG}。
    \end{enumerate}

\subsubsection{配置 HMAC 密钥和密钥功能}
        \begin{enumerate}
            \item 将密钥功能 \begin{math}m\end{math} 写入寄存器 \hyperref[regdesc:HMACSETPARAPURPOSEREG]{HMAC\_SET\_PARA\_PURPOSE\_REG}。表 \ref{tab:hmac-key-purpose} 列出了 \begin{math}m\end{math} 对应的数值，可参见章节 \ref{sec:hmac-config-function}。
            \item 将数值 \textcolor{red}{n} 写入寄存器 \hyperref[regdesc:HMACSETPARAKEYREG]{HMAC\_SET\_PARA\_KEY\_REG}，选择 eFuse 存储器中的 KEY\textcolor{red}{n} 作为本次计算的密钥（\textcolor{red}{n} 的取值范围为 0 \verb+~+ 5\tagged{ESP32-P4-latest}{，7}），可参见章节 \ref{sec:hmac-config-purposes}。
            \item 将 1 写入寄存器 \hyperref[regdesc:HMACSETPARAFINISHREG]{HMAC\_SET\_PARA\_FINISH\_REG}，完成配置。
            \item 读取寄存器 \hyperref[regdesc:HMACQUERYERRORREG]{HMAC\_QUERY\_ERROR\_REG}。如果返回值为 1，表明选取的密钥块与配置的密钥功能不匹配，本次计算任务中止。如果返回值为 0，表明选取的密钥块与配置的密钥功能匹配，可以执行计算流程。
            \item 如果 \hyperref[regdesc:HMACSETPARAPURPOSEREG]{HMAC\_SET\_PARA\_PURPOSE\_REG} 的数值不为 8，表明 HMAC 模块当前处于下行模式，请参考章节 \ref{sec:hmac-downstream-process} 进行下一步配置。如果数值为 8，表明 HMAC 模块当前处于上行模式，请参考章节 \ref{sec:hmac-upstream-process} 进行下一步配置。
        \end{enumerate}

\subsubsection{下行模式流程}
\label{sec:hmac-downstream-process}
        \begin{enumerate}
            \item 轮询状态寄存器 \hyperref[regdesc:HMACQUERYBUSYREG]{HMAC\_QUERY\_BUSY\_REG}，当读取到该寄存器的值为 0 时，继续下一步骤。
            \item 如需清除计算结果，并确保后续不再使用相关硬件（JTAG 或 RSA\_DS），用户可将 1 写入寄存器 \hyperref[regdesc:HMACSETINVALIDATEJTAGREG]{HMAC\_SET\_INVALIDATE\_JTAG\_REG}，清除 JTAG 密钥生成的结果；也可将 1 写入寄存器 \hyperref[regdesc:HMACSETINVALIDATEDSREG]{HMAC\_SET\_INVALIDATE\_DS\_REG}，清除 RSA\_DS\_KEY 生成的结果。此后，如需重启 JTAG 或 RSA\_DS 外设，需重新进行 HMAC 调用流程。
        \end{enumerate}

\subsubsection{上行模式流程}
\label{sec:hmac-upstream-process}

\begin{enumerate}

    \item 上行模式下传输数据：
    \begin{enumerate}
        \item 轮询状态寄存器 \hyperref[regdesc:HMACQUERYBUSYREG]{HMAC\_QUERY\_BUSY\_REG}，当读取到该寄存器的值为 0 时，继续下一步骤。
        \item 用户根据章节 \ref{sec:hmac-padding} 的描述对数据进行 SHA 附加填充并获取数据块 Block\_\textcolor{red}{n} ($\textcolor{red}{n}>=0$)
        \item 根据数据块 \textcolor{red}{n} 的数量，选择其中一种操作：
        \begin{itemize}
            \item 如果 Block\_\textcolor{red}{n} 是唯一一个数据块，同时 Block\_\textcolor{red}{1} 已经包含了所有的填充位：
            \begin{enumerate}
                \item 将 512 位的数据块 Block\_\textcolor{red}{n} 写入寄存器  \hyperref[regdesc:HMACWRMESSAGENREG]
                {HMAC\_WR\_MESSAGE\_\regindex{n}\_REG (\regindex{n}: 0-15)} 中。随后，将 1 写入寄存器 \hyperref[regdesc:HMACSETMESSAGEONEREG]{HMAC\_SET\_MESSAGE\_ONE\_REG} 中，HMAC 模块将计算该数据块。
                \item 轮询状态寄存器 \hyperref[regdesc:HMACQUERYBUSYREG]{HMAC\_QUERY\_BUSY\_REG}，当读取到该寄存器的值为 0 时，继续下一步骤。
                \item 写 1 到寄存器 \hyperref[regdesc:HMACONEBLOCKREG]{HMAC\_ONE\_BLOCK\_REG} 中，完成数据传输。
            \end{enumerate}
            \item 如果 Block\_\textcolor{red}{n} 是最后一个填充数据块：
            \begin{enumerate}
                \item 将 512 位的数据块 Block\_\textcolor{red}{0} 写入寄存器 \hyperref[regdesc:HMACWRMESSAGENREG]
                {HMAC\_WR\_MESSAGE\_\regindex{n}\_REG (\regindex{n}: 0-15)} 中。 随后，将 1 写入寄存器 \hyperref[regdesc:HMACSETMESSAGEONEREG]{HMAC\_SET\_MESSAGE\_ONE\_REG} 中, HMAC 模块将计算该数据块。
                \item 轮询状态寄存器 \hyperref[regdesc:HMACQUERYBUSYREG]{HMAC\_QUERY\_BUSY\_REG}，当读取到该寄存器的值为 0 时，继续下一步骤。
                \item 自动完成数据传输。
            \end{enumerate}
            \item 如果 Block\_\textcolor{red}{n} 是倒数第二个填充数据块：
            \begin{enumerate}
                \item 将 512 位的数据块 Block\_\textcolor{red}{n} 写入寄存器 \hyperref[regdesc:HMACWRMESSAGENREG]
                {HMAC\_WR\_MESSAGE\_\regindex{n}\_REG (\regindex{n}: 0-15)} 中。随后，将 1 写入寄存器 \hyperref[regdesc:HMACSETMESSAGEONEREG]{HMAC\_SET\_MESSAGE\_ONE\_REG} 中，HMAC 模块将计算该数据块。
                \item 轮询状态寄存器 \hyperref[regdesc:HMACQUERYBUSYREG]{HMAC\_QUERY\_BUSY\_REG}，当读取到该寄存器的值为 0 时，继续下一步骤。
                \item 写 1 到寄存器 \hyperref[regdesc:HMACSETMESSAGEPADREG]{HMAC\_SET\_MESSAGE\_PAD\_REG} 中，随后跳转到步骤 c 继续数据传输。
            \end{enumerate}

            \item 如果 Block\_\textcolor{red}{n} 既不是最后一个也不是倒数第二个数据块：
            \begin{enumerate}
                \item 将 512 位的数据块 Block\_\textcolor{red}{n} 写入寄存器 \hyperref[regdesc:HMACWRMESSAGENREG]
                {HMAC\_WR\_MESSAGE\_\regindex{n}\_REG (\regindex{n}: 0-15)} 中。随后，将 1 写入寄存器 \hyperref[regdesc:HMACSETMESSAGEONEREG]{HMAC\_SET\_MESSAGE\_ONE\_REG} 中，HMAC 模块将计算该数据块。
                \item 轮询状态寄存器 \hyperref[regdesc:HMACQUERYBUSYREG]{HMAC\_QUERY\_BUSY\_REG}，当读取到该寄存器的值为 0 时，继续下一步骤。
                \item 写 1 到寄存器 \hyperref[regdesc:HMACSETMESSAGEINGREG]{HMAC\_SET\_MESSAGE\_ING\_REG} 中，随后跳转到步骤 c 继续数据传输。
            \end{enumerate}
        \end{itemize}
    \end{enumerate}

    \item 读取上行模式下的结果 hash 值：
    \begin{enumerate}
        \item 轮询状态寄存器 \hyperref[regdesc:HMACQUERYBUSYREG]{HMAC\_QUERY\_BUSY\_REG}，当读取到该寄存器的值为 0 时，继续下一步骤。
        \item 从寄存器 \hyperref[regdesc:HMACRDRESULTNREG]{HMAC\_RD\_RESULT\_\regindex{n}\_REG (\regindex{n}: 0-7)} 中读取 hash 结果数值。
        \item 将 1 写入寄存器 \hyperref[regdesc:HMACSETRESULTFINISHREG]{HMAC\_SET\_RESULT\_FINISH\_REG}，结束当次计算，并同时清除计算结果。
        \item 上行模式下的操作完成。
    \end{enumerate}
	\end{enumerate}


\begin{tiplisting}
\label{hmac_tips_2}
    RSA\_DS 外设和 HMAC 模块可直接调用或在内部使用 SHA 加速器，但两者不能同时共享硬件资源。因此在 HMAC 模块运行过程中，SHA 模块无法被 CPU 和 RSA\_DS 外设调用。
\end{tiplisting}

\section{HMAC 算法细节}

\subsection{附加填充位}
\label{sec:hmac-padding}

HMAC 模块中采用 SHA-256 作为加密 hash 算法。该算法中，若待输入数据的总位数不是 512 的倍数，用户须在软件中应用 SHA-256 附加填充算法。SHA-256 附加填充算法与 \href{https://nvlpubs.nist.gov/nistpubs/FIPS/NIST.FIPS.180-4.pdf}{FIPS PUB 180-4} 中章节 \textit{Padding the Message} 所描述相同。下行模式下，用户无需输入数据或进行数据填充，HMAC 模块默认使用 32 字节 0x00 方式重启 JTAG，以及 32-byte 0xff 方式生成 RSA\_DS 外设的 AES 密钥。

为方便阅读，以下将简要介绍数据填充的步骤。如图 \ref{fig:hmac-padding} 所示，假设待处理数据长度为 \begin{math}m\end{math} 个位，填充步骤如下：
\begin{enumerate}
    \item 在待处理数据末尾附加一个 “1”。
    \item 附加 \begin{math}k\end{math} 个 “0”。其中，\begin{math}k\end{math} 为满足 $m + 1 + k {\equiv} 448 (mod 512)$ 的最小非负数。
    \item 附加一个 64 位的整数值作为二进制块。该二进制块的内容为待填充数据作为一个大端二进制整数值 \begin{math}m\end{math} 的长度。
\end{enumerate}

\begin{figure}[ht!]
    \centering
    \fbox{\includegraphics[width=0.7\textwidth]{\modulefiles/figures/HMAC-padding__CN}}
    \caption{HMAC 附加填充位示意图}
    \label{fig:hmac-padding}
\end{figure}

上行模式下，若待填充数据总位数是 512 的整数倍，则用户可通过将 1 写入 \hyperref[regdesc:HMACSETMESSAGEENDREG]{HMAC\_SET\_MESSGAE\_END\_REG} 配置由硬件完成 SHA 附加填充操作，也可通过将 1 写入 \hyperref[regdesc:HMACSETMESSAGEPADREG]{HMAC\_SET\_MESSAGE\_PAD\_REG} 自行完成填充操作。若待填充数据总位数不是 512 的整数倍，则用户只能自行完成 SHA 附加填充操作。数据填充操作完成后，用户应参照章节 \ref{sec:hmac-process-detailed} 完成后续配置。

\subsection{HMAC 算法结构}
HMAC 模块中应用的算法结构示意图如 \ref{fig:hmac-structure} 所示。这是 RFC 2104 中描述的标准 HMAC 算法。

\begin{figure}[ht!]
    \centering
    \fbox{\includegraphics[width=0.65\textwidth]{\modulefiles/figures/HMAC-structure__CN}}
    \caption{HMAC 结构示意图}
    \label{fig:hmac-structure}
\end{figure}

图 \ref{fig:hmac-structure} 中，
\begin{enumerate}
    \item ipad 是由 64 个 0x36 字节组成的 512 位数据块。
    \item opad 是由 64 个 0x5c 字节组成的 512 位数据块。
\end{enumerate}

首先，HMAC 模块在 256 位的密钥 K 的位序列后附加上 256 位的 0 序列，得到 512 位的 K$_0$。再对 K$_0$ 和 ipad 进行异或运算，得到 512 位的 S1。将总位数为 512 倍数的待输入数据附加到 512 位的 S1 数值后，使用 SHA-256 加密算法计算得到 256 位的 H1。

HMAC 模块通过对 K$_0$ 和 opad 进行异或运算得到 S2，将 256 位的 hash 计算结果附加到 512 位的 S2 数值后，得到 768 位长度的序列，使用章节 \ref{sec:hmac-padding} 中描述的 SHA 附加填充算法将该序列填充成 1024 位的序列，最后使用 SHA-256 加密算法计算得到的最终 hash 结果（256 位）。

% >>> If your module has no registers, delete sample text
%       for Sections Register Summary and Registers below <<<

%%%%%%%%%%%%%%%%%%%%%%%%
%%% Register Summary %%%
%%%%%%%%%%%%%%%%%%%%%%%%

% >>> Update the sample text below and insert
%     the info on register summary into the subfile <<<

\newpage
\hypertarget{hmac-reg-summ}{}
\section{寄存器列表}
\label{sec:hmac-reg-summ}

本小节的所有地址均为相对于 HAMC 加速器基地址的地址偏移量（相对地址），具体基地址请见章节 \ref{mod:sysmem} \textit{\nameref{mod:sysmem}} 中的表 \ref{tab:sysmem-base-address}。

请查看章节 \hyperref[glossary-access-types]{\textit{寄存器的访问类型}}，了解“访问”列缩写的含义。

\subfile{\modulefiles/28-HMAC-reg-summ__CN}


%%%%%%%%%%%%%%%%%
%%% Registers %%%
%%%%%%%%%%%%%%%%%

% >>> Update the sample text below and insert
%      the info on registers into the subfile <<<

\newpage
\section{寄存器}

本小节的所有地址均为相对于 HAMC 加速器基地址的地址偏移量（相对地址），具体基地址请见章节 \ref{mod:sysmem} \textit{\nameref{mod:sysmem}} 中的表 \ref{tab:sysmem-base-address}。

关于保留 (reserved) 域的处理，请查看章节 \hyperref[programming-reserved-field]{如何配置寄存器的保留域}。

\subfile{\modulefiles/28-HMAC-reg__CN}

\end{document}
